% =============================================================================
% cap6 §6.1 Síntese
% =============================================================================
\section{Síntese}
\label{sec:conclusao-sintese}

A monografia abordou o problema de manejo de colônia de gatos no Campus~2 da USP São Carlos do ponto de vista da engenharia de sistemas, propondo um \textbf{pipeline de visão computacional em quatro fases} --- triagem por movimento, detecção genérica de fauna, classificação de espécie e re-identificação individual --- como instrumento de redução de fricção operacional para o projeto \emph{Gatosdoc2}. O Capítulo~\ref{cap:hardware} caracterizou os pontos de instalação e fixou a tipologia única de \emph{kit} de câmeras IP comerciais domésticas alimentadas por baterias, com comunicação Wi-Fi entre câmera e \emph{gateway} local. O Capítulo~\ref{cap:pipeline} desenhou o pipeline em servidor único, sem inferência embarcada na borda, e fixou \modeloMD, \modeloSN, \modeloPetFace{} e o método de partes do corpo de \citeauthor{akbar2025bodypart} sobre \modeloWildlife{} como componentes principais, recortando os pontos de \emph{tradeoff} privacidade-LGPD, custo de inferência e cobertura taxonômica. O Capítulo~\ref{cap:metodologia} formalizou o protocolo de avaliação restrito a datasets públicos, com baseline mínimo por fase e teste integrador sobre janela sintética. O Capítulo~\ref{cap:resultados} reportou as estimativas conceituais correspondentes, com elementos fictícios devidamente sinalizados, e discutiu as implicações arquiteturais.
